\chapter{知识的迁移与适应}
\label{chap:ket-transfer-adaptation}

花卉识别系统已经能够区分一批常见物种。把它交给另一个植物园时，新的要求却未必是继续识别这批物种：园内可能出现来源数据中没有的花，也可能要求判断叶片病害，或在整幅照片中标出每朵花的位置。已有模型中的颜色、纹理和形状表示可能仍然有用，但原来的输出规则不一定回答新的问题。需要解决的是：哪些来源经验能够支撑目标能力，目标还须提供什么证据，以及形成这种能力需要多少样本和计算。

第~\ref{chap:ket-update-maintenance}章关注系统改变以后怎样维持、修正和继续更新已有能力，本章则考察来源知识怎样服务于一个新的目标。二者可以同时发生。例如，学习新物种属于本章的适应问题；把新物种加入长期服务并保持旧物种识别能力，还须满足前章的维护要求。判断一种迁移方案时，应把目标语义、可用信息和计算权限放在一起考虑，再选择合适的准备方式与目标端操作。少样本和零样本描述信息条件，\term{元学习}（meta-learning）则从历史任务中学习怎样适应其他任务，也就是用过去的学习经验改进后来的学习过程。它们不是相互排斥的目标任务类别。

\section{可迁移的目标任务}
\label{sec:ket-transfer-targets}

一次迁移并不由“换了一个数据集”充分定义。设来源任务的输入、输出空间为$X_S,Y_S$，联合分布为$\mathcal P_S$；目标端对应为$X_T,Y_T,\mathcal P_T$。目标风险写作
\begin{equation}
  R_T(f)=\mathbb E_{(\bm{x},y)\sim\mathcal P_T}
  \bigl[\ell_T(f(\bm{x}),y)\bigr].
  \label{eq:ket-transfer-target-risk}
\end{equation}
这里的损失$\ell_T$表达目标任务所认可的错误，未必与来源损失相同。输入分布、输出语义和评价准则分别变化，都可能使原模型不再适用。目标任务还包括训练时能看到的信息、允许执行的适应操作，以及未来待预测样本的范围。这些条件确定以后，“迁移得好”才有可以检验的含义。

\subsection{预测目标与输出语义的确定}
\label{sec:ket-transfer-output-semantics}

新增一种花时，若仍要求每张图输出一个物种名称，变化主要在类别集合与分类边界。判断病害则改变了标签含义：同一物种可以健康，也可以感染某种病害；把物种编号改成病害编号，并不能使原分类头自动获得病害知识。定位花朵的变化更大，输出需要包含位置、数量以及每个对象的类别。即使把输出向量加长，模型仍须知道哪些坐标表示同一个对象、怎样处理不同数量的花、如何将预测与标注框配对。分割需要逐像素标签，生成需要对输出序列建模；它们都要求与新语义相符的预测结构、监督与损失。

类别集合的关系只在类别含义已经对齐后才有意义。令$C_S,C_T$分别表示源、目标中实际可能出现的语义类别，$C_{\cap}=C_S\cap C_T$为共享部分。类别名称相同不保证定义相同，例如来源按物种划分，目标按园艺品种划分时，应先明确映射，不能直接比较编号。表~\ref{tab:ket-transfer-class-relations}固定“单幅图像分类”以及常见的无监督领域自适应协议：训练时有带标签源样本和无标签目标样本，目标标签只用于锁定方案后的评价。类别关系与监督协议是两个维度；若论文允许少量目标标签、源数据不可访问或采用源自由适应，必须另行注明，不能只沿用表中的设定名称。

\begin{table*}[htbp]
  \centering
  \small
  \setlength{\tabcolsep}{3pt}
  \begin{tabularx}{\linewidth}{p{19mm}p{25mm}p{29mm}XX}
    \toprule
    设定 & 类别关系 & 代表性训练访问 & 目标端判断 & 额外风险 \\
    \midrule
    完全共享 & $C_T=C_S$ & 带标签源样本；无标签目标样本 & 在共同类别中分类 & 拍摄条件变化仍可破坏原分类器 \\
    部分领域自适应 & $C_T\subsetneq C_S$ & 同左；训练时不揭示目标标签 & 迁移目标实际具有的类别 & 源中特有类别可能被错误地对齐到目标 \\
    开放集领域自适应 & 典型设定为$C_S\subsetneq C_T$ & 同左；目标私有类不提供类别名 & 共享类分类，目标私有类拒识为未知 & 未知不能被强迫匹配到已知类 \\
    通用领域自适应 & 两侧均可有私有类 & 同左；共享范围预先未知 & 共享类分类，其余目标样本输出未知 & 共享范围与未知比例都需要估计 \\
    \bottomrule
  \end{tabularx}
  \caption{源目标类别关系、代表性监督协议及预测责任。四行采用带标签源样本与无标签目标样本的常见协议；其他监督或访问条件须另列。开放集一行采用源类别全部属于目标已知类别的常用设定；这里的未知输出不是各个目标私有类别的名称。}
  \label{tab:ket-transfer-class-relations}
\end{table*}

来源类别较多、目标类别较少时，\term{部分领域自适应}（partial domain adaptation）要从来源中选出对目标有用的类别，避免多余内容造成负迁移。部分对抗领域自适应（Partial Adversarial Domain Adaptation, PADA）利用目标样本对源类别的预测分数估计类别权重，降低源中特有类别在分类和领域对抗训练中的作用。它没有把所有源样本同等地向目标对齐，但权重来自当前模型，早期预测错误也可能压低真正共享类别的权重\scite{ket-cao2018partial}

目标存在源中没有的类别时，强行对齐全部目标输入会把未知物种推向已知物种。\term{开放集领域自适应}（open set domain adaptation）为模型增加“未知”的选项；Saito等的方案通过分类器与特征生成器的对抗过程，让目标样本可以被对齐为已知，也可以与已知分离\scite{ket-saito2018openset}\term{通用领域自适应}（universal domain adaptation）进一步不预设源目标标签集的包含关系，即允许两边都具有不属于对方的类别；模型利用样本可迁移性判断应迁移哪些共享内容，并把目标私有类别统一标为未知。这里的“通用”是标签集合关系上的放宽，不表示对任意新任务都能成功\scite{ket-you2019universal}

未知判断有三个不同层次。把一张花卉图像标为“已知类别之外”，只完成了\term{未知拒识}（unknown rejection），即承认现有类别目录不足。把拒识样本分成若干稳定群组，还需要发现它们内部的结构，且群组可能反映背景或拍摄角度而非物种。要学会一个新类别，则须确定其语义、获得足以识别它的证据，并在新的样本上验证预测规则。聚类后经专家命名、补充代表性标注、重新训练分类器，是一种可能的流程；开放集评价中的一个“未知”标签并没有替我们完成后面这些工作。

例如，来源类别为$\{\text{菊花},\text{月季},\text{鸢尾}\}$，目标为$\{\text{月季},\text{鸢尾},\text{新种甲}\}$。系统应迁移月季和鸢尾的相关知识，排除菊花对目标分类的不当干扰；没有新种甲的可用语义或样本时，只能争取将它拒识为未知。若目标随后提供新种甲的描述和标注，才有依据建立这个新类别。若还要求菊花在旧植物园继续可识别，则需要同时满足第~\ref{sec:ket-update-retention-validation}节的旧能力保留要求。

任务变化与分布变化并不互斥。类别完全相同而相机、季节变化时，目标语义可以不变，但$\mathcal P_S$与$\mathcal P_T$不同；拍摄条件相近而标签从物种改为病害时，输入边缘分布甚至可以相同，预测目标仍然改变。后者不能仅靠匹配输入分布解决。第~\ref{chap:transfer-learning-theory}章说明了缺少跨域标签联系时的不可识别边界。需要分布校正时，调用第~\ref{sec:ket-update-weighting}节的加权机制；观测空间改变时，按式~\eqref{eq:ket-update-crossmodal}等约束建立有依据的对应，并检查新表示是否保留目标所需的信息。“领域适应”也常用来称呼这里的领域自适应，本书统一使用后者。

\subsection{目标信息与适应权限的确定}
\label{sec:ket-transfer-evidence}

面对新种甲和新种乙，目标植物园可以只提供文字描述，也可以提供每类几张带标签图像，还可以允许访问一批没有标签的当地照片。这三种条件提供的信息不同。描述给出类别语义，标注图像展示语义怎样落实到具体观测，无标签图像揭示目标输入的分布与局部结构。即使三个方案都只做一次前向预测，也不能把它们视为使用了相同证据。

\term{少样本学习}（few-shot learning）指目标端只有少量可用于学习的标注样本。分类任务常写成$N$类、每类$K$例，即$N$-way $K$-shot，此时支持样本共$NK$个；还应报告每类查询数量、类别不平衡、重复图像与增强方式。\term{零样本学习}（zero-shot learning）则必须附带“什么没有见过”的说明。没有目标类别的训练图像，不等于没有类别名称和属性；语言模型没有目标任务示范，也不等于预训练文本从未描述过该任务。图像零样本与指令零示范使用的缺失条件不同，不能只凭“零”字比较知识需求。

为控制信息差异，可以为每个新物种任务约定一个证据包：类别描述$D_T$、带标签支持集$S_T$、允许使用的无标签集$U_T$以及已有交互历史$H_T$。这些符号表示信息来源，不要求每种方案都必须使用全部信息。比较时除固定证据包外，还要记录来源模型是否有相同预训练数据，描述是否经过专家整理，检索库是否含有目标示范，以及模型选择是否利用额外标签。

访问权限决定证据怎样进入计算。\term{归纳式预测}（inductive prediction）在适应后独立处理每个新输入，其他查询输入不影响当前结果；\term{直推式预测}（亦称转导式预测，transductive prediction）允许联合利用整批未标注查询输入，例如根据它们的聚类结构调整分类边界。直推式方法可以合法地使用这些输入，但不能使用其隐藏答案，也不能把整批输入带来的信息增益归因于优化器本身。逐样本部署只有一个输入到来时，整批查询方法还可能无法原样执行。

另一些权限也会改变问题：能否更新骨干参数，能否保存跨样本状态，能否在不确定时向专家查询，是否允许在新请求到来前离线计算。主动查询会增加新证据，其获取与校验属于第~\ref{part:efficient-knowledge-use}部分的知识获取主题；本章只在预算中记录允许查询多少次、何时获得反馈。不能让系统在预测一个查询后读取其正确标签，再把更新后的结果当作原来无反馈条件下的测试结果。

因此，“快速适应”应落实为带条件的比较：给定标注数$n$、更新步数$T$和计算预算$B$，比较目标风险，或比较达到指定风险所需的代价。类别描述整理、示范检索、提示选择和验证集调参都可能比一次梯度更新更昂贵。只报告目标端反向传播次数，会漏掉不更新参数方案的主要费用。

\subsection{历史任务与目标任务关系的确定}
\label{sec:ket-transfer-task-splits}

历史任务只有与未来任务共享某些规律，才能为未来学习提供依据。设一个任务$\mathcal T_k$包含输入输出语义、联合分布$\mathcal P_k$、损失以及信息权限，任务从分布$\Pi$中抽取。$\Pi$不是单个图像的分布：先抽取“在某个地区区分哪几种花”这样的学习问题，再从该问题中抽取图像，才是两层抽样。历史任务可能共同依赖花瓣形状，也可能只在背景上相似；选择历史任务时需要说明预计保留的规律是什么。

对每个任务，将用于适应的\term{支持集}（support set）$S_k$与用于检查适应结果的\term{查询集}（query set）$Q_k$分开。支持集告诉学习者当前任务怎样定义，查询集检验由这些证据得到的规则能否处理其他样本。二者在同一任务内可以包含相同类别，但不应含有同一照片、同一原图的增强副本，或按评价单位应共同隔离的同一植株记录。任务之间还应分为元训练、元验证和最终测试任务；检验未见类别迁移时，类别也要在这三部分之间隔离。

把准备好的共享知识记为$\bm{\phi}$，把基于支持集产生任务预测器的过程记为$A_{\bm{\phi}}(S_k)$。跨任务准备的一个典型目标为
\begin{equation}
  \min_{\bm{\phi}}
  \mathbb E_{\mathcal T_k\sim\Pi}
  \mathbb E_{S_k,Q_k\mid\mathcal T_k}
  \bigl[\hat R_{Q_k}(A_{\bm{\phi}}(S_k))\bigr],
  \label{eq:ket-transfer-meta-objective}
\end{equation}
其中$\hat R_{Q_k}$是查询集上的平均损失。元训练查询集在任务内扮演检验角色，却在任务间扮演训练信号，因此不是最终独立测试集。模型结构、步长和停止时机在元验证任务上选择；最终目标任务只用允许的支持证据适应，其查询答案留到方案冻结后评价。图~\ref{fig:ket-transfer-data-boundaries}把这两种反馈区分开来。

\begin{figure*}[htbp]
  \centering
  % Editable schematic, designed at 166 mm width; no measured data.
% Solid arrows: evidence use / prediction. Dashed arrows: training feedback.
\begin{tikzpicture}[
  x=1mm,y=1mm,
  font=\figurefont\footnotesize,
  text=BookInk,
  >={Latex[length=2mm,width=1.4mm]},
  evidence/.style={draw=BookBlue,fill=FigureInput!10,line width=0.8pt,
    text width=24mm,minimum height=12mm,align=center,inner sep=2mm},
  model/.style={draw=BookTeal,fill=FigureModel!10,line width=0.8pt,
    text width=26mm,minimum height=12mm,align=center,inner sep=2mm},
  score/.style={draw=BookAmber,fill=FigureLoss!10,line width=0.8pt,
    text width=28mm,minimum height=12mm,align=center,inner sep=2mm},
  flow/.style={draw=BookInk,line width=0.9pt,-{Latex},rounded corners=1mm},
  feedback/.style={draw=BookAmber,line width=0.9pt,dashed,-{Latex},
    rounded corners=1mm}
]
  \node[anchor=west,font=\figurefont\small] at (2,15) {历史任务：允许跨任务训练反馈};
  \node[evidence] (support-one) at (16,0) {$S_1$：月季、鸢尾\\支持图像与标签};
  \node[model] (adapt-one) at (59,0) {用共享准备$\bm\phi$\\形成任务预测器1};
  \node[score] (query-one) at (101,0) {$Q_1$：其他图像\\查询预测与损失};

  \node[evidence] (support-two) at (16,-29) {$S_2$：菊花、百合\\支持图像与标签};
  \node[model] (adapt-two) at (59,-29) {用共享准备$\bm\phi$\\形成任务预测器2};
  \node[score] (query-two) at (101,-29) {$Q_2$：其他图像\\查询预测与损失};
  \node[model,text width=25mm,minimum height=19mm] (prepare) at (150,-14.5)
    {跨任务更新\\共享准备$\bm\phi$};

  \draw[flow] (support-one.east) -- (adapt-one.west);
  \draw[flow] (adapt-one.east) -- (query-one.west);
  \draw[flow] (support-two.east) -- (adapt-two.west);
  \draw[flow] (adapt-two.east) -- (query-two.west);
  \draw[feedback] (query-one.east) -- (123,0) |- (prepare.west);
  \draw[feedback] (query-two.east) -- (123,-29) |- (prepare.west);
  \node[anchor=west,text=BookMuted] at (2,-43)
    {任务1、2均从同一轮准备出发；更新后进入下一轮历史任务};

  \draw[draw=BookRule,line width=0.8pt] (2,-50) -- (165,-50);
  \node[anchor=west,font=\figurefont\small] at (2,-59)
    {留出目标：方案冻结后，只用支持证据适应};

  \node[evidence] (support-target) at (16,-76)
    {$S_T$：新种甲、乙\\允许的图像与标签};
  \node[model] (adapt-target) at (59,-76)
    {目标任务预测器\\由$S_T$形成};
  \node[model,text width=23mm,fill=FigureOutput!8] (predict-target) at (103,-76)
    {$Q_T$输入\\目标预测};
  \node[score,text width=24mm] (score-target) at (151,-76)
    {$Q_T$隐藏答案\\最终评分};
  \draw[flow] (support-target.east) -- (adapt-target.west);
  \draw[flow] (adapt-target.east) -- (predict-target.west);
  \draw[flow] (predict-target.east) -- (score-target.west);
  % Frozen transfer is a one-way route outside the historical task lanes.
  \draw[flow] (prepare.south) -- (150,-47) -- (167,-47)
    -- (167,-64) -- (59,-64) -- (adapt-target.north);
  \node[anchor=east,fill=white,inner sep=1mm,text=BookTeal]
    at (148,-64) {冻结准备结果};
  \node[anchor=west,text=BookMuted] at (2,-92)
    {目标评分没有返回适应或模型选择的箭头};
\end{tikzpicture}

  \caption{任务内适应与任务间学习的数据边界。两个历史花卉任务都利用共享准备和支持集，通过参数适应或条件化形成任务预测器；查询损失以虚线反馈到共享准备，实线表示使用证据与预测计算。留出的目标任务只接收冻结后的准备结果，其查询答案仅用于最终评分，不返回适应或模型选择。类别与任务划分均为教学构造，留出任务也可以同时包含领域变化。}
  \label{fig:ket-transfer-data-boundaries}
\end{figure*}

Baxter的归纳偏置学习框架解释了为何多个相关任务可以帮助选择可复用的假设空间：每个候选空间都应允许任务内部拟合不同的函数，而跨任务经验用来比较哪个空间经常包含好的任务解。单个任务的样本增加，有助于辨认该任务中的预测规则；独立任务数量增加，才为“这种限制是否适用于其他任务”提供新证据。两种样本量不能互相完全替代，而且候选假设空间族的容量也必须受控，否则准备阶段仍可能记住有限历史任务\scite{ket-baxter2000-bias}

未见类别不等于未见领域。在相同拍摄设备上识别另一组物种，与从温室照片迁移到野外红外影像，考验的是不同的覆盖。Meta-Dataset通过多种数据来源以及变化的类别数和支持样本数，让评价不只依赖固定的五分类任务；跨领域少样本评测还要求将目标领域整体留出\scite{ket-triantafillou2019metadataset}解释差异时，应分别考察类别是否新、领域是否新、每类样本是否变少，以及类别之间是否更难区分，而不是把它们统称为“更难的任务”。

有限历史任务可以通过合成方式扩展，但合成不能凭空证明未来覆盖。任务插值方法在历史任务之间混合中间特征；标签语义共享时可以同步插值标签，不共享时则配对类别并为合成类别重新赋予任务内标签。支持集与查询集分别构造，不能把查询答案混入支持证据。生成的“介于两个历史任务之间”的问题是否有助于真实新地区、新物种，需要在独立目标任务上验证\scite{ket-yao2021interpolation}

还有一种更隐蔽的失败：模型在元训练中表现很好，却没有使用支持集。若所有任务都把月季标为1，且输入暴露了任务身份，模型可以直接记住类别编号。随机交换任务内标签、替换支持集或移除任务身份，可以检验预测是否真正依赖适应证据。Yin等进一步研究了这种记忆问题，以信息约束限制模型从参数或查询输入直接记住训练任务的捷径，促使它使用支持信息；这并不意味着任意打乱标签都能解决所有任务的记忆问题\scite{ket-yin2019memorization}通用的数据隔离见第~\ref{sec:ket-evaluation-splits}节，任务层重复与不确定性分析见第~\ref{sec:ket-evaluation-repeats}节。

\section{可迁移知识的准备}
\label{sec:ket-transfer-preparation}

来源经验可以在目标学习中承担不同功能。表示使样本可比较，语义联系使新类别有可识别的含义，初始化与先验限制参数的起点和可取范围，更新器规定怎样利用反馈，条件模型则把任务证据转成当前预测所需的状态。这些功能可以共存在一个模型中。已有预训练模型已经提供足够表示时，不必为了使用少量目标样本再做一次元训练；若未来任务很多且结构稳定，专门准备适应过程才可能值得。

\subsection{可迁移表示的准备}
\label{sec:ket-transfer-representation}

假设来源编码器把图像映射为$\bm z=h_{\bm{\theta}}(\bm x)\in\mathbb R^d$。若目标物种在这个空间中已经容易分开，少量标注只需确定类别对应的区域；若病害表现为来源任务刻意忽略的叶片细节，再复杂的线性分类头也无法从已经丢失的信息中恢复病害。判断表示的作用，可以固定来源数据、骨干、图像增强和目标划分，比较冻结表示后的线性分类、近邻以及类别均值，而不是同时更换表示和适应算法。

普通监督预训练就可以形成这种表示。Chen等在一致实现下比较少样本分类方法，并用跨领域设置检查普通微调基线；Tian等研究冻结表示后拟合线性分类器的基线，还考察了自蒸馏对表示的影响。这些工作说明，表示质量和实现条件可能解释相当一部分基准差异，但不能推出所有新任务都只需复用特征。涉及自蒸馏时，其额外准备成本以及第~\ref{sec:ket-integration-student-training}节的教师学生训练过程也应计入比较\scite{ket-chen2019closer,ket-tian2020embedding}

如果准备阶段明确知道目标将只有少量支持样本，可以直接训练表示，使其适合这种使用方式。\term{原型网络}（Prototypical Networks）为每个任务内类别计算一个表示均值，再按查询与均值的距离分类。令$S_{k,c}$为任务$k$中标签为$c$的支持子集，原型与预测概率为
\begin{align}
  \bm c_{k,c}
  &=\frac{1}{|S_{k,c}|}
    \sum_{(\bm x_i,y_i)\in S_{k,c}}h_{\bm{\theta}}(\bm x_i),
    \label{eq:ket-transfer-prototype}\\
  p_{\bm{\theta}}(y=c\mid\bm x,S_k)
  &=\frac{\exp\{-\|h_{\bm{\theta}}(\bm x)-\bm c_{k,c}\|_2^2/\tau\}}
    {\sum_{c'=1}^{N}\exp\{-\|h_{\bm{\theta}}(\bm x)-\bm c_{k,c'}\|_2^2/\tau\}}.
    \label{eq:ket-transfer-prototype-probability}
\end{align}
其中$\tau>0$控制距离转成概率后的尖锐程度，标准形式可取$\tau=1$。均值并非对任何比较规则都天然合适：在平方欧氏距离下，令$\sum_i\|\bm z_i-\bm c\|_2^2$对$\bm c$的导数为零，得到$\bm c=|S_{k,c}|^{-1}\sum_i\bm z_i$。因此它确实最小化同类支持点到代表点的平方距离，但这一准则没有保证多峰类别能被单个点充分表示\scite{ket-snell2017-prototypical}

一次原型网络训练会从历史类别中抽出$N$类，为每类抽$K$个支持样本，再从剩余图像抽查询样本。编码器处理两组输入，支持表示形成原型，查询按式~\eqref{eq:ket-transfer-prototype-probability}得到类别概率。任务损失为
\begin{equation}
  L_k(\bm{\theta})
  =-\frac{1}{|Q_k|}\sum_{(\bm x,y)\in Q_k}
  \log p_{\bm{\theta}}(y\mid\bm x,S_k).
  \label{eq:ket-transfer-episodic-loss}
\end{equation}
梯度同时经过查询表示和支持原型回传到编码器。原型不是在元训练前计算一次就固定：编码器变化后，它们必须用当前支持表示重算。不断更换历史任务重复这个过程，保存的是表示参数和比较约定，而不是把训练任务原型当作未来类别。

\term{匹配网络}（Matching Networks）保留每个支持样本的作用，不先将同类样本压成均值。令$f_{\bm{\theta}}(\bm x,S_k)$和$g_{\bm{\theta}}(\bm x_i,S_k)$分别为查询、支持表示，把两者的相似度记为$s_i(\bm x,S_k)$，则
\begin{align}
  s_i(\bm x,S_k)
  &=\operatorname{sim}\bigl(f_{\bm{\theta}}(\bm x,S_k),
       g_{\bm{\theta}}(\bm x_i,S_k)\bigr), \nonumber\\
  a_i(\bm x,S_k)
  &=\frac{\exp\{s_i(\bm x,S_k)/\tau\}}
  {\sum_{j=1}^{|S_k|}\exp\{s_j(\bm x,S_k)/\tau\}}, \nonumber\\
  p_{\bm{\theta}}(y=c\mid\bm x,S_k)
  &=\sum_{i=1}^{|S_k|}a_i(\bm x,S_k)\,\mathbf 1\{y_i=c\}.
  \label{eq:ket-transfer-matching}
\end{align}
这就是按相似程度给支持标签加权投票。原论文使用余弦比较，还研究完整上下文表示：支持编码器通过双向循环网络读取支持集合的一个排列，查询编码器通过多次注意力读取支持表示。由此，同一图像的表示可以随当前候选集合变化；这种实现额外引入了支持排列和循环计算，不能把它简化成所有样本相互独立的固定特征提取\scite{ket-vinyals2016matching}

匹配网络也用式~\eqref{eq:ket-transfer-episodic-loss}训练，只是把预测概率换成式~\eqref{eq:ket-transfer-matching}。逐例匹配保留局部模式，但查询计算和任务存储随支持样本数增加；原型网络降低了这两项代价，却要求均值足以代表类别。在每类只有一例、两者使用相同编码器和相同比较分数时，两种规则一致；若一个使用余弦、另一个使用平方距离，或仅一个使用完整上下文编码器，则不能仅凭“一例”宣称等价。目标端冻结表示后仍要建立参照，这项实际计算将在第~\ref{sec:ket-transfer-references}节展开。

\subsection{可迁移语义联系的准备}
\label{sec:ket-transfer-semantic-preparation}

没有新物种的训练图像时，支持样本的均值无从计算。不过，目标可以提供“细长花瓣、叶缘有锯齿”等描述，来源经验则可以教会系统识别这些属性。关键是把“知道类别应具有什么属性”和“能够从图像看出这些属性”分开：前者来自语义描述，后者来自带相应监督的来源观测。只有两者连接起来，文字才可能支持图像预测。

以二值属性为例，设$\bm a_c\in\{0,1\}^m$是类别$c$的属性描述。对来源图像，用其类别属性或逐图标注训练$m$个属性预测器$q_j(\bm x)\approx p(a_j=1\mid\bm x)$。\term{直接属性预测}（direct attribute prediction, DAP）通过这些预测器估计图像属性，再把属性与目标描述匹配；属性预测器的输出维度固定，不因目标类别增加而增加。若用类别描述给每张来源图赋相同属性标签，就引入了“同类对象共有这些属性”的简化，对类内变异和不可见属性需要谨慎\scite{ket-lampert2009attributes}

怎样由属性预测得到类别判断？在DAP中，类别$c$确定属性向量$\bm a_c$，即$p(\bm a\mid y=c)=\mathbf 1\{\bm a=\bm a_c\}$；模型还假定给定属性后，类别判断不再依赖图像。记$\pi_c=p(y=c)$，在相应属性向量的先验概率为正时，贝叶斯公式与对属性的边缘化给出
\begin{equation}
  \begin{aligned}
    p(y=c\mid\bm a)
      &=\frac{\pi_c\,\mathbf 1\{\bm a=\bm a_c\}}{p(\bm a)},\\
    p(y=c\mid\bm x)
      &=\sum_{\bm a\in\{0,1\}^m}
          p(y=c\mid\bm a)\,p(\bm a\mid\bm x)\\
      &=\pi_c\,\frac{p(\bm a_c\mid\bm x)}{p(\bm a_c)}.
  \end{aligned}
  \label{eq:ket-transfer-attribute-link}
\end{equation}
类别确定属性，使求和中只留下$\bm a=\bm a_c$这一项；到这一步还没有把属性分布拆成乘积。为用单个属性预测器计算比值，另作两项近似：给定图像时各属性条件独立，以及属性先验可以因子化。令$p_j=p(a_j=1)$，则分别有
\[
  \begin{aligned}
    p(\bm a_c\mid\bm x)
      &\approx\prod_{j=1}^{m}
        q_j(\bm x)^{a_{c,j}}[1-q_j(\bm x)]^{1-a_{c,j}},\\
    p(\bm a_c)
      &\approx\prod_{j=1}^{m}
        p_j^{a_{c,j}}(1-p_j)^{1-a_{c,j}}.
  \end{aligned}
\]
第一项近似并不推出第二项：属性即使在固定图像条件下独立，对不同图像混合以后仍可能相关。把这两项近似代入式~\eqref{eq:ket-transfer-attribute-link}，得到用于比较目标类别的分数
\begin{equation}
  s_c(\bm x)=\pi_c
  \prod_{j=1}^{m}
  \frac{q_j(\bm x)^{a_{c,j}}
        [1-q_j(\bm x)]^{1-a_{c,j}}}
       {p_j^{a_{c,j}}(1-p_j)^{1-a_{c,j}}},
  \label{eq:ket-transfer-attribute-score}
\end{equation}
分子衡量图像与描述的符合程度，分母按属性在来源中的常见程度作校正。Lampert等按来源类别属性表逐列取均值估计属性先验，而不是按来源图像频数加权。本节沿用逐类等权的口径，并另加预先指定的对称伪计数$\delta>0$作平滑：
\[
  \hat p_j=\frac{\sum_{c\in C_S}a_{c,j}}{|C_S|},
  \qquad
  p_j=\frac{\sum_{c\in C_S}a_{c,j}+\delta}{|C_S|+2\delta}.
\]
这里$C_S$非空，平滑保证$p_j\in(0,1)$；某类的图像更多，不会使该类在此估计中获得更大权重。目标类别先验在本节取均匀值$\pi_c=1/|C_T|$，不利用待评价查询的标签估计类别频率。实际计算通常累加对数以免连乘下溢。这些分布近似与来源估计未必符合真实目标分布，分数归一化后也不能自动视为经过校准的真实类别后验。

Lampert等还比较了间接属性预测：先估计来源类别概率，再通过来源类别的属性表计算
\[
  p(a_j=1\mid\bm x)
  =\sum_{c\in C_S}a_{c,j}\,p(y=c\mid\bm x).
\]
它复用来源分类器，却也受来源类别覆盖限制。无论采用哪条路径，若新种甲和新种乙具有完全相同的属性描述，且类别先验相同，式~\eqref{eq:ket-transfer-attribute-score}就无法区分二者。若唯一不同的是叶缘特征，而照片没有拍到叶片，则描述虽然不同，当前观测仍不足。更多语义词条不能替代可观察且可识别的区分证据。

显式属性需要事先选定属性词表，连续图文表示提供了另一种联系。\term{对比语言图像预训练}（Contrastive Language--Image Pre-training, CLIP）用配对图像和文本训练两个编码器，使真实配对的表示相近。来源文本因而不只提供固定类别编号，还提供自然语言中的视觉描述。预训练的完整对比学习过程属于第~\ref{part:efficient-knowledge-use}部分的知识获取主题；这里需要保存的是图像编码器、文本编码器、共同表示空间及其尺度约定。目标类别描述经过文本编码器，便能作为图像比较的参照\scite{radford2021}

连续语义表示能够容纳多种表述，但某个向量维度不必对应一个可解释属性，相似度也不保证能精确区分否定、计数或细粒度形态。类别名称、属性和解释都是目标知识，描述由谁提供、是否经过目标验证集选择、预训练是否已经覆盖目标物种，都需要记录。严格的“来源未见类”实验必须控制来源类别；大规模网页预训练的零样本评价通常只能承诺没有用目标训练集进行有监督拟合，不能据此宣称模型从未接触相关概念。

\subsection{参数初始化与任务先验的准备}
\label{sec:ket-transfer-initialization-prior}

表示与语义参照约束了怎样解释输入，另一种准备方式则约束任务模型怎样开始学习。初始化$\bm{\theta}$是一个参数点；任务先验$p(\bm w_k\mid\bm\eta)$是任务参数可能取值的分布，$\bm\eta$为它的共享超参数。前者可以配合梯度更新，后者可以配合后验推断。一个先验的均值可以被用作初始化，但仅有均值并没有指定先验的方差、相关性或多峰结构。

\subsubsection{参数初始化的学习}
\label{sec:ket-transfer-maml}

普通预训练寻找在来源数据上直接预测得好的参数，新任务却可能只允许做几次更新。因此，也可以把“更新后是否预测得好”直接作为准备目标。\term{模型无关元学习}（model-agnostic meta-learning, MAML）学习一个共享起点，让它经过支持集上的少量梯度步骤后，在同任务查询集上表现良好。这里的模型无关是指适用于可通过相应梯度过程训练的参数化模型，不是无需假设或无需选定模型结构\scite{finn2017}

记$L_{S_k}(\bm w)$、$L_{Q_k}(\bm w)$为任务$k$的支持、查询平均损失。对一个含$M$个历史任务的批次，一步适应的内层和外层分别为
\begin{align}
  \bm w_k^{(0)}&=\bm{\theta}, \nonumber\\
  \bm w_k^{(1)}&=\bm{\theta}
       -\alpha\nabla_{\bm{\theta}}L_{S_k}(\bm{\theta}),
       \label{eq:ket-transfer-maml-inner}\\
  J(\bm{\theta})&=\frac{1}{M}\sum_{k=1}^{M}
       L_{Q_k}(\bm w_k^{(1)}(\bm{\theta})), \nonumber\\
  \bm{\theta}&\leftarrow\bm{\theta}-\beta\nabla_{\bm{\theta}}J(\bm{\theta}).
       \label{eq:ket-transfer-maml-outer}
\end{align}
$\alpha$是任务内步长，$\beta$是跨任务步长。每个任务都从当前同一个$\bm{\theta}$独立出发，而不是让任务$k+1$接着任务$k$的适应参数继续训练。外层更新修改的是这个共享起点，不是把各任务的最终分类器直接保存为一个分类器。

外层梯度须经过适应步骤回传。令$\bm H_k=\nabla^2L_{S_k}(\bm{\theta})$，按链式法则有
\begin{align}
  \frac{\partial\bm w_k^{(1)}}{\partial\bm{\theta}}
     &=\bm I-\alpha\bm H_k, \nonumber\\
  \nabla_{\bm{\theta}}J(\bm{\theta})
     &=\frac{1}{M}\sum_{k=1}^{M}
       (\bm I-\alpha\bm H_k)^{\top}
       \nabla_{\bm w}L_{Q_k}(\bm w_k^{(1)}).
       \label{eq:ket-transfer-maml-chain}
\end{align}
查询梯度指出适应后参数应怎样改变，前面的雅可比矩阵则把这个要求换算成对初始参数的要求。对足够光滑的损失，Hessian对称；实际自动微分可计算Hessian与向量的乘积，不必显式存储完整平方大小的矩阵。

若执行$T$步内层更新，令
\[
  \bm J_k^{(0)}=\bm I,\qquad
  \bm J_k^{(t+1)}
  =\bigl[\bm I-\alpha_t\nabla^2L_{S_k}(\bm w_k^{(t)})\bigr]
    \bm J_k^{(t)}.
\]
则第$k$项的元梯度为$(\bm J_k^{(T)})^\top\nabla L_{Q_k}(\bm w_k^{(T)})$。这个递推固定了矩阵相乘的顺序；用反向自动微分实现时沿更新轨迹反向传播。普通混合任务训练则直接最小化$\sum_kL_{S_k}(\bm{\theta})$或合并数据损失，没有“用支持集形成任务解，再在查询集上评价”的中间映射。两者的目标不同，即使骨干和优化器相同，也不应只比较更新前的分类准确率。

算法~\ref{alg:ket-transfer-maml}把参数重置和数据边界放到同一个执行过程中。目标使用阶段只执行其中的任务内更新，跨任务更新属于准备费用。

\begin{algorithm}[MAML的准备与目标适应]
  \label{alg:ket-transfer-maml}
  \begin{algorithmic}[1]
    \Require 历史训练任务、独立元验证任务；内层步数$T$与步长$\alpha_t$；外层步长$\beta$
    \Ensure 冻结的共享初始化及给定目标支持集的适应模型
    \State 初始化共享参数$\bm{\theta}$
    \While{训练预算尚未耗尽，且未满足元验证停止条件}
      \State 抽取$M$个历史任务，各自构造不重叠的$S_k,Q_k$
      \For{$k=1,\ldots,M$}
        \State 复制共享起点：$\bm w_k^{(0)}\gets\bm{\theta}$
        \For{$t=0,\ldots,T-1$}
          \State $\bm w_k^{(t+1)}\gets\bm w_k^{(t)}
            -\alpha_t\nabla L_{S_k}(\bm w_k^{(t)})$
        \EndFor
        \State 用$\bm w_k^{(T)}$计算查询损失$L_{Q_k}$
      \EndFor
      \State 沿内层更新轨迹求导，用平均查询梯度更新$\bm{\theta}$
    \EndWhile
    \State 按元验证结果选定并冻结初始化、步长与步数
    \State 新目标从该初始化出发，仅在$S_T$上执行约定的$T$步更新
    \State 对$Q_T$输入产生预测；隐藏答案只进入独立评分
  \end{algorithmic}
\end{algorithm}

\begin{example}[同一组分类任务中的起点与适应后损失]
  \label{ex:ket-transfer-maml-numeric}
  为了可手算，设花卉图像已压成一维特征$z$，用$wz$预测二值标签$y\in\{-1,1\}$的分数，以$\operatorname{sign}(wz)$分类，以平方损失训练。以下数值是教学构造，不是MAML的实验结果。任务1的支持记录为$(1,1),(-1,-1)$，查询记录为$(2,1),(-2,-1)$；任务2的支持记录为$(2,1),(-2,-1)$，查询记录为$(1,1),(-1,-1)$。每条记录来自独立样本，支持与查询的特征幅度分布有意设为不同。

  对每组采用$\frac{1}{2n}\sum_i(wz_i-y_i)^2$，得到
  \[
    L_{S_1}(w)=\tfrac12(w-1)^2,\quad
    L_{S_2}(w)=\tfrac12(2w-1)^2,
  \]
  \[
    L_{Q_1}(w)=\tfrac12(2w-1)^2,\quad
    L_{Q_2}(w)=\tfrac12(w-1)^2.
  \]
  直接混合两个支持损失的最优起点为$w_{\mathrm{pre}}=3/5$。令内层步长$\alpha=1/4$，两个任务的一步参数分别为$w_1^{(1)}=3w/4+1/4$和$w_2^{(1)}=1/2$。代入查询损失，元目标为
  \[
    J(w)=\tfrac14(3w/2-1/2)^2+\tfrac1{16},
  \]
  最优元起点为$w_{\mathrm{meta}}=1/3$。在$w=3/5$处，$\mathrm dJ/\mathrm dw=0.3$，其中任务2的内层雅可比$1-\alpha\cdot4=0$，所以它对元梯度的贡献为零，可以在求和中略去；不能把这个雅可比因子当作1，直接用适应后的查询梯度代替对起点的梯度。

  两个起点更新前的平均查询损失分别为$0.05$与$5/36\approx0.1389$，普通起点更小；相同的一步更新后却分别为$0.1025$与$0.0625$，元起点更小。此例中分数的符号均正确，分类准确率不能区分这些起点，差别体现在训练所用的分数损失。它说明评价对象和更新预算会改变比较结论，不证明元学习在真实花卉任务中必然占优。
\end{example}

完整MAML需要保存或重算内层计算图。\term{一阶MAML}（first-order MAML, FOMAML）把适应雅可比近似为单位矩阵，用适应后查询梯度更新起点，省去经过支持梯度的二阶求导。它仍然需要采样历史任务、执行内层更新和计算查询反馈；“一阶”减少的是求导路径，不是删除元训练。

Reptile采用另一种一阶过程：从共享起点独立训练每个历史任务若干步，得到$\bm w_k^{(T)}$，然后更新
\begin{equation}
  \bm{\theta}\leftarrow\bm{\theta}
   +\epsilon\frac1M\sum_{k=1}^{M}
        \bigl(\bm w_k^{(T)}-\bm{\theta}\bigr).
  \label{eq:ket-transfer-reptile}
\end{equation}
它不要求像MAML那样另取查询集计算外层梯度；任务内小批次的多步训练本身产生更新方向。只有一步普通SGD时，该方向退化为混合任务梯度；多步时，后续梯度是在前面更新后的参数处计算，包含任务内不同小批次相互作用的信息。因此，Reptile不是“把MAML的Hessian删掉”的另一个名字。独立元验证与最终目标测试仍必须保留，其不使用显式查询梯度不构成测试泄漏的例外\scite{ket-nichol2018reptile}

两步展开可以看出这种差异。把同一任务两个小批次在共同起点的梯度记为$\bm g_1,\bm g_2$，第二批损失的Hessian记为$\bm H_2$。在足够光滑且步长较小的条件下，
\[
  \bm w^{(2)}-\bm{\theta}
    =-\alpha(\bm g_1+\bm g_2)
      +\alpha^2\bm H_2\bm g_1+O(\alpha^3).
\]
第二批梯度在$\bm{\theta}-\alpha\bm g_1$处求值，因此出现$\bm H_2\bm g_1$这一相互作用项。Reptile实现只计算一阶梯度，并不显式计算这个Hessian；泰勒展开是在解释多步更新包含什么信息，不是额外增加的算法步骤。

适应的收益究竟来自表示变化，还是来自一个已经很好的表示，也可以拆开检查。将模型参数分成骨干$\bm{\theta}_b$和任务头$\bm{\theta}_h$，\term{几乎无内循环}（Almost No Inner Loop, ANIL）在每个任务的内层只更新任务头，外层仍训练骨干与头的初始化。查询损失既直接依赖骨干产生的查询表示，也通过支持表示影响头的适应过程；因此，精确ANIL并不等于把骨干在整个元训练期间冻结。原研究在MiniImageNet等少样本任务和相应网络中观察到强特征复用作用，但新目标需要此前未保留的局部信息时，头部更新仍可能不足\scite{ket-raghu2019anil}

\subsubsection{任务参数先验的学习}
\label{sec:ket-transfer-prior}

只有两张病害图像时，多组参数都可能与它们相符。选择一个参数点并不能表达这种不确定性。层次先验把历史任务中反复出现的规律写成任务参数分布，再让新证据修正该分布。设$\bm\eta$为跨任务超参数，$\bm w_k$为任务参数，给定各任务输入后，层次模型分解为
\begin{equation}
  p(\bm\eta,\{\bm w_k,S_k\}_{k=1}^{M})
   =p(\bm\eta)\prod_{k=1}^{M}
      p(\bm w_k\mid\bm\eta)
      \prod_{(\bm x_i,y_i)\in S_k}
      p(y_i\mid\bm x_i,\bm w_k).
  \label{eq:ket-transfer-hierarchical-model}
\end{equation}
式中把输入作为条件略去；给定共享超参数后，任务参数相互独立，给定任务参数后，任务内样本按所设似然生成。历史任务可以用积分掉$\bm w_k$后的边缘似然估计共享超参数：
\begin{equation}
  \hat{\bm\eta}
  =\argmax_{\bm\eta}\sum_{k=1}^{M}
     \log\int p(S_k\mid\bm w_k)\,
          p(\bm w_k\mid\bm\eta)\,\mathrm d\bm w_k.
  \label{eq:ket-transfer-empirical-bayes}
\end{equation}
这称为经验贝叶斯：历史数据用来估计先验分布的参数，而不是假定先验完全由人工给定。另一种做法是保留$\bm\eta$的后验。二者都要求说明所用近似和任务关系，不能把一个正则项的系数随意命名为“学到的先验”。此处的历史支持集用于估计先验，准备方案仍由独立历史验证任务选择，最终目标查询答案不进入这一估计。

在目标端，若准备阶段已经固定$\hat{\bm\eta}$，支持证据产生
\begin{align}
  p(\bm w_T\mid S_T,\hat{\bm\eta})
  &\propto p(S_T\mid\bm w_T)p(\bm w_T\mid\hat{\bm\eta}), \nonumber\\
  \hat{\bm w}_T
  &=\argmin_{\bm w}\left[
       -\sum_{(\bm x_i,y_i)\in S_T}\log p(y_i\mid\bm x_i,\bm w)
       -\log p(\bm w\mid\hat{\bm\eta})\right].
  \label{eq:ket-transfer-map}
\end{align}
第一行是后验分布，第二行仅是最大后验点估计。使用完整后验预测时，还须计算$\int p(y\mid\bm x,\bm w)p(\bm w\mid S_T,\hat{\bm\eta})\,\mathrm d\bm w$，不能用一个点的输出冒充参数不确定性已经被积分掉。

一个可解析的例子说明先验强度。设$\bm w_k\sim\mathcal N(\bm\mu,\sigma_w^2\bm I)$，连续目标满足$y_i=\bm z_i^\top\bm w_k+\varepsilon_i$，$\varepsilon_i\sim\mathcal N(0,\sigma_y^2)$。这是对花卉某个连续测量的简化回归模型，并非直接把类别编号当连续值。历史任务估计共享中心$\bm\mu$和任务差异尺度$\sigma_w^2$；对新任务，把支持特征按行组成$\bm Z$，观测组成$\bm y$，由似然与先验的二次项相加得到
\begin{align}
  \bm\Sigma_T
    &=\left(\sigma_y^{-2}\bm Z^\top\bm Z
                     +\sigma_w^{-2}\bm I\right)^{-1}, \nonumber\\
  \bm m_T
    &=\bm\Sigma_T\left(\sigma_y^{-2}\bm Z^\top\bm y
                     +\sigma_w^{-2}\bm\mu\right).
  \label{eq:ket-transfer-gaussian-posterior}
\end{align}
这个模型的后验为$\mathcal N(\bm m_T,\bm\Sigma_T)$，均值也等于最大后验估计。支持数据不仅改变中心，还改变各方向的方差；只有初始化$\bm\mu$并不能提供$\bm\Sigma_T$。

\begin{example}[先验怎样改变少样本估计]
  \label{ex:ket-transfer-prior-numeric}
  取一维模型$z_i=1$，设来源任务给出的先验为$w_T\sim\mathcal N(0,1)$，观测噪声方差为1，两个教学观测为$2,4$。不使用先验时估计为样本均值3；后验精度为$1+2=3$，均值为$(0+2+4)/3=2$，方差为$1/3$。若先验方差减为$1/4$，均值变为$6/(4+2)=1$，方差为$1/6$；若先验方差增为4，均值为$6/(1/4+2)=8/3$，方差为$4/9$。强先验减少不确定性，却也更容易压低真正偏离历史中心的新任务参数。
\end{example}

初始化与先验之间确有联系，但联系取决于推断过程。Grant等把MAML解释为层次模型中的一种经验贝叶斯近似：用少量梯度步骤形成任务参数点估计，再近似适应后的预测评价\scite{ket-grant2018hierarchical}

一维二次损失可以说明早停何时对应正的先验精度。设负对数似然的有效二次项为$L(w)=h(w-\hat w)^2/2$，曲率$h>0$，从$w^{(0)}=\mu$以固定步长$\alpha$更新$T\geq1$步。记$a=(1-\alpha h)^T$，递推得到$w^{(T)}=(1-a)\hat w+a\mu$。另一方面，
\[
  \argmin_w\left\{L(w)+\frac{\lambda}{2}(w-\mu)^2\right\}
    =\frac{h\hat w+\lambda\mu}{h+\lambda}.
\]
当$0<\alpha h<1$时，$0<a<1$，取$\lambda=ha/(1-a)>0$便得到相同参数点，对应高斯先验$\mathcal N(\mu,\lambda^{-1})$。例如$h=1,\alpha=1/2,T=1$给出$a=1/2,\lambda=1$。这只是在指定损失尺度下的点估计对应，精度依赖曲率、步长和步数；换成非线性网络以后，有限步梯度既不保证找到全局后验众数，也没有恢复完整后验。多个已知任务共同先验的结构设计留给第~\ref{sec:ket-sharing-prior}节，本节的目标是把历史任务学到的先验交给未来任务使用。

\subsection{参数更新规则的准备}
\label{sec:ket-transfer-learned-optimizer}

同一个起点配合不同步长、动量和梯度缩放，会产生不同适应轨迹。除了学习起点，还可以学习“看到这样的梯度以后怎样移动”。设任务模型参数为$\bm w_k^{(t)}$，更新器参数为$\bm\psi$，内部状态为$\bm s_k^{(t)}$。先按保存的模型起点与新任务头初始化规则建立$\bm w_k^{(0)}$，再由初始化函数$\operatorname{Init}_{\bm\psi}$建立更新器状态，之后才开始递推：
\begin{align}
  \bm s_k^{(0)}
    &=\operatorname{Init}_{\bm\psi}(\bm w_k^{(0)}), \nonumber\\
  \bm g_k^{(t)}
    &=\nabla_{\bm w}L_{S_k}^{(t)}(\bm w_k^{(t)}), \nonumber\\
  (\Delta\bm w_k^{(t)},\bm s_k^{(t+1)})
    &=U_{\bm\psi}(\bm g_k^{(t)},\bm w_k^{(t)},
             L_{S_k}^{(t)},\bm s_k^{(t)}), \nonumber\\
  \bm w_k^{(t+1)}
    &=\bm w_k^{(t)}+\Delta\bm w_k^{(t)}.
  \label{eq:ket-transfer-learned-update}
\end{align}
这里$L_{S_k}^{(t)}$允许取支持集中的不同小批次；输入哪些量由具体更新器决定。$\bm w_k$随目标反馈改变，$\bm\psi$则在历史任务上训练并在目标适应时固定。把两者混称为模型参数，会掩盖究竟迁移了预测能力还是学习过程。

Andrychowicz等用循环网络参数化更新器。为避免让更新器规模随任务模型参数数目平方增长，同一个小型LSTM逐坐标共享权重，而每个坐标保留各自的状态；它读取该坐标梯度及历史，输出坐标增量。梯度幅度和符号需要适当预处理，使不同量级的优化任务能被同一网络处理。历史梯度因此可以影响当前更新，作用类似于可学习的动量与步长规则，但不能由此推断它具有手工优化器在某个凸问题上的收敛保证\scite{ket-andrychowicz2016-learning}

在逐坐标实现中，$\operatorname{Init}_{\bm\psi}$为每个坐标建立独立状态，可置零，也可复制准备阶段学得的初态；所选规则须与$\bm\psi$一同保存。每个新任务或独立重复运行都重新建立模型起点并调用$\operatorname{Init}_{\bm\psi}$，同一任务内则在支持小批次之间沿用状态，更新器只读取许可的支持反馈。上一任务的末态不自动传给下一任务；跨任务保留状态须另设连续适应协议，最终查询标签始终不进入任务内更新。

为训练更新器，把式~\eqref{eq:ket-transfer-learned-update}展开$T$步，对历史任务的最终查询损失或多个时刻的查询损失求和：
\begin{equation}
  J(\bm\psi)=\frac1M\sum_{k=1}^{M}
       \sum_{t=1}^{T}\omega_t L_{Q_k}(\bm w_k^{(t)}(\bm\psi)),
       \qquad \omega_t\geq0.
  \label{eq:ket-transfer-optimizer-objective}
\end{equation}
少样本任务使用独立查询数据，可鼓励更新器学到适应后的泛化；纯函数优化也可直接评价同一目标函数值，两种训练目的应区别。反向传播沿模型参数和更新器状态的轨迹更新$\bm\psi$。截断轨迹会改变可传递的长期信息，停止对输入梯度再求导会省掉一部分高阶项；这些都是计算近似，不是对原目标精确求导的别名。

Ravi与Larochelle将少样本分类器的参数与门控循环状态联系起来，用学习到的保留门和更新门控制参数变化。其核心可写成$\bm w^{(t+1)}=\bm f^{(t)}\odot\bm w^{(t)}-\bm i^{(t)}\odot\bm g^{(t)}$：$\bm f$决定保留多少当前参数，$\bm i$决定各坐标的更新幅度，门由当前损失、梯度、参数和历史状态形成。初始状态也可训练，从而联合学习起点与更新规则；查询标签只在元训练外层评价，不交给任务内更新器\scite{ket-ravi2017optimization}

VeLO把准备任务扩展到更大范围的神经网络优化问题，用任务和训练进度等信息调节学习到的优化行为。它说明更新器也能成为可复用的训练资产，但不能把论文中无需逐任务搜索学习率的使用方式解释为全过程没有调参或计算成本。任务模型更大、训练轨迹更长、损失类型或梯度尺度改变时，均需重新检查其外推表现，尤其要测量在更长训练中是否稳定，而不仅看短轨迹的终点\scite{ket-metz2022velo}

保存更新器需要$|\bm\psi|$个参数，执行时还可能为任务模型的每个坐标保留循环状态。若任务模型有$P$个参数，每坐标状态宽度为$h$，仅这些状态就需$O(Ph)$存储；更新器每一步的额外前向计算也应加到普通梯度计算上。某个任务少调几次学习率可能节省使用费用，却未必抵消跨任务训练、更新器存储和每步执行开销。

\subsection{任务条件化机制的准备}
\label{sec:ket-transfer-conditioning-preparation}

任务证据不一定要通过梯度写入参数。系统也可以先读支持样本，形成当前任务的摘要，再根据摘要回答查询。\term{任务条件化}（task conditioning）就是让当前预测显式依赖任务证据；共享参数规定怎样读取证据，任务状态则随这次证据变化。与第~\ref{sec:ket-transfer-learned-optimizer}节不同，这里可以不产生一套经优化得到的任务权重。

\term{条件神经过程}（Conditional Neural Process, CNP）给出一种清楚的实现。对支持记录$(\bm x_i,y_i)$，编码器$e_{\bm\phi}$产生$\bm r_i\in\mathbb R^{d_r}$，将其平均后交给条件预测器：
\begin{align}
  \bm r_k&=\frac1{|S_k|}\sum_{(\bm x_i,y_i)\in S_k}
                   e_{\bm\phi}(\bm x_i,y_i), \nonumber\\
  p_{\bm\phi}(y\mid\bm x,S_k)
      &=p_{\bm\phi}(y\mid\bm x,\bm r_k).
  \label{eq:ket-transfer-cnp}
\end{align}
平均操作使$\bm r_k$不依赖支持样本排列；均值也可能丢失样本数和多峰细节，需要时应将计数或其他统计量显式加入接口。回归解码器可以输出均值和正方差，分类解码器可以输出类别概率。网络参数$\bm\phi$在历史任务上由条件负对数似然训练，目标端只计算新的$\bm r_T$。

原始CNP训练从一组观测中抽取上下文，预测目标可包含上下文点以及未观察点。若沿用本章的支持查询评价，泛化损失应在与支持样本分离的$Q_k$上报告；上下文重构可以是辅助训练信号，不能替代独立查询评价。在编码器宽度固定、每点计算固定的简单均值实现中，处理$n$个支持点和$q$个查询点约为$O(n+q)$次点级网络计算，支持摘要还可以缓存。

固定支持集时，基本因子化解码器把多个查询输出建模为条件独立，未表示它们的完整相关性。换用不同支持集时，这些条件预测也不保证都能由同一个随机函数先验经贝叶斯条件化得到；分别输出预测均值与方差，并没有自动保证这种共同的概率模型解释。预测方差是否校准还须另行检验，不能由模型“给出方差”推出\scite{ket-garnelo2018cnp}

上下文示例也可以通过序列模型读取。MetaICL在历史任务中抽取若干支持输入答案对与一个独立查询，把支持记录和查询输入拼成序列，用查询答案的负对数概率更新模型。准备目标可写为
\begin{equation}
  \min_{\bm\phi}
  \mathbb E_k\left[
    -\log p_{\bm\phi}
       (y\mid D_k,(\bm x_1,y_1),\ldots,(\bm x_n,y_n),\bm x)
  \right],
  \label{eq:ket-transfer-metaicl}
\end{equation}
其中末尾$(\bm x,y)$来自该历史任务的查询集，查询答案用于损失，不作为预测它自身的输入。$D_k$表示可选的任务说明，是本书为覆盖“说明加示范”的扩展接口；原始MetaICL的核心输入是任务内示范与查询，并不要求每个任务另有自然语言说明。训练完成后，模型在未见任务中读取新示例而不更新权重，这就是显式训练上下文适应能力的准备路径\scite{ket-min2022metaicl}

另一条路径是直接使用已经预训练的语言模型。GPT-3的少样本使用把示例放入提示，不对每个目标任务执行梯度微调；其准备过程主要是大规模语言建模，而非按式~\eqref{eq:ket-transfer-metaicl}显式构造的任务训练。因此，“模型输入里有示例”不能证明它受过元学习训练，也不能证明输出变化只来自那些示例。要记录的是模型实际保存了什么条件化能力、覆盖了哪些任务，以及目标应按什么格式提供证据\scite{brown2020}

准备结束时，应保存表示、语义编码器、先验超参数、更新器及其初始化规则和附加模块，并把预训练或元训练计入式~\eqref{eq:ket-reuse-amortization}的$C_{\mathrm{prep}}$。沿用该式的完整账目，若将来源选择费用并入本章的$C_{\mathrm{adapt},k}$，以$C_{\mathrm{predict},k}$表示单次预测费用，$K$个实际目标任务的计算费用简记为
\begin{equation}
  C_{\mathrm{total}}
    =C_{\mathrm{prep}}
     +\sum_{k=1}^{K}
       \bigl(C_{\mathrm{adapt},k}+q_kC_{\mathrm{predict},k}\bigr),
  \label{eq:ket-transfer-total-cost}
\end{equation}
其中$q_k$为任务$k$的预测次数，检索和状态管理按实际发生阶段归入目标端费用。准备成本的摊销、失败尝试与无迁移对照沿用第~\ref{sec:ket-reuse-resource-benefits}节，不再另设口径。元训练的效果仍须在历史覆盖之外的目标上检验。

\section{目标任务的适应}
\label{sec:ket-transfer-adaptation}

准备结果交付给目标端以后，适应可以改变三种对象：显式比较的预测参照，条件模型读取的上下文，以及通过优化得到的模型参数。参照和上下文都可能包含带标签样本，却不必更新网络权重；参数优化则把目标反馈写入可训练变量。三种操作允许组合，例如先用文字描述建立类别参照，再用少量图像拟合一个校正头，不必让每个任务顺次执行全部操作。

\subsection{预测参照的构建}
\label{sec:ket-transfer-references}

沿用新种甲与新种乙，假定当前任务只要求在这两个指定的未见类别中分类。这个候选集只含未见类，属于传统零样本学习（ZSL）的评价条件；若测试时已见类与未见类同时出现，则属于广义零样本学习（GZSL），候选集合应改为$C_S\cup C_T$，并分别报告已见类与未见类表现及其折中。传统ZSL上的归一化概率和阈值不能直接外推到GZSL，因为来源类别通常具有更强的训练优势。没有目标图像时，目标须提供足以区分新类别的描述。对DAP，可以直接把属性向量代入式~\eqref{eq:ket-transfer-attribute-score}；对CLIP类模型，可以将文本描述编码并归一化为$\bm t_c$，将查询图像编码并归一化为$\bm z$，计算
\begin{equation}
  p(y=c\mid\bm x,D_T)
  =\frac{\exp(\bm z^\top\bm t_c/\tau)}
         {\sum_{c'=1}^{N}\exp(\bm z^\top\bm t_{c'}/\tau)}.
  \label{eq:ket-transfer-text-reference}
\end{equation}
参照的内容来自目标描述，比较能力来自来源图文联系。固定模型不代表参照固定：加入一个新描述会改变候选集合，softmax的归一化概率也随之改变。多个描述模板可以先分别编码，再按照预先约定的方式平均并重新归一化；选择模板所用的验证信息与费用不能隐藏。

\begin{example}[描述、支持样本与原型的比较计算]
  \label{ex:ket-transfer-reference-numeric}
  以下属性和向量都是为演示构造的，不是实际植物分类依据，也不是任何模型的实测输出。设两个教学属性为“花瓣细长”和“叶缘有锯齿”，新种甲的描述为$(1,0)$，新种乙为$(0,1)$。查询属性预测为$(0.8,0.3)$，两类先验均为$1/2$。假设来源属性表中每列恰有一半类别取值为1，则前述逐类均值及对称平滑都给出属性先验$1/2$。DAP中公共因子消去后，两个分数正比于$0.8(1-0.3)=0.56$和$(1-0.8)0.3=0.06$，归一化分数为$28/31\approx0.9032$与$3/31\approx0.0968$。

  若共同语义空间中的归一化文本参照为$\bm t_{\text{甲}}=(1,0)$、$\bm t_{\text{乙}}=(0,1)$，查询图像表示为$\bm z=(0.8,0.6)$，取$\tau=0.2$，则两类分数为4和3，式~\eqref{eq:ket-transfer-text-reference}给出甲类概率$e^4/(e^4+e^3)\approx0.7311$。两个示例都选择甲，但概率来自不同表示与假设，数值不能直接当作两种方法的可靠性比较。

  现在提供每类两张支持图像，另用一个冻结视觉编码器得到甲类表示$(0,0),(0,2)$，乙类表示$(2,0),(2,2)$，查询为$(0.5,1)$。原型为$\bm c_{\text{甲}}=(0,1)$、$\bm c_{\text{乙}}=(2,1)$，查询到原型的平方距离为$0.25,2.25$。取$\tau=1$，甲类概率为$1/(1+e^{-2})\approx0.8808$。若用相同的负平方距离做逐例匹配，四个支持点的平方距离为$1.25,1.25,3.25,3.25$，同类权重相加得到同一概率。这里的相等来自这个对称构造，而不是两种方法在多样本时恒等。
\end{example}

原型的便宜有明确代价。考虑另一组一维教学表示：甲类支持点为$-2,2$，乙类为$0.5,0.7$，查询为2。两类原型分别为0与0.6，平方距离为4与1.96，单原型判为乙；使用负平方距离的逐例匹配时，甲的未归一化总权重为$1+e^{-16}$，乙为$e^{-2.25}+e^{-1.69}$，结果判为甲。甲类由两个相距很远的模式组成，平均后落在没有支持样本的位置。可以增加类内原型或保留逐例参照，但每类只有极少样本时，又未必足以稳定估计多个模式。

来源参照不能不加判断地保留。部分重叠任务需要识别$C_{\cap}$，避免把来源私有类别作为目标必须选择的答案；开放集任务则需要允许不属于任何已知参照的输入退出。一个简单规则是先计算已知类别分数，再采用
\begin{equation}
  \hat y(\bm x)=
  \begin{cases}
    \argmax_{c\in C_{\mathrm{known}}}p(c\mid\bm x),
       & \max_c p(c\mid\bm x)\geq\gamma,\\
    \mathrm{unknown}, & \text{否则}.
  \end{cases}
  \label{eq:ket-transfer-rejection}
\end{equation}
阈值$\gamma$须依据独立验证数据、可接受的误拒率或成本约束选择，不能用最终未知测试集调到最好再报告。softmax强制概率和为1，即使所有已知类别都不合适，也可能出现很高的最大概率；距离阈值、未知验证类别或专门的拒识模型可以提供额外依据，但都带来假设和信息需求。只用已知类验证可以约束已知类误拒，不能单独保证任意未知类的检出率。

如果目标从分类变为检测，可以在局部候选区域使用类别参照，但还需要产生候选区域、预测位置并处理重复对象；分割需要保留空间分辨率并生成像素级输出；生成需要把参照转成解码条件。分类原型本身没有这些输出结构。相对地，拟合一个线性分类器虽然也基于冻结特征，但它用优化改变了任务权重，应归入第~\ref{sec:ket-transfer-parameter-adaptation}节，而不与直接求均值统称为“不训练”。

\subsection{任务上下文的组织}
\label{sec:ket-transfer-context}

同一张花卉照片可以回答“属于哪种花”，也可以回答“是否需要复查叶片病害”。一个条件模型需要知道当前问题、示范怎样给出答案，以及哪些部分尚待预测。目标端的\term{上下文学习}（in-context learning, ICL）指模型根据当前输入中的示例改变本次任务行为；若执行过程没有梯度更新，这种改变发生在激活、注意力读取或任务状态中，不自动成为跨会话保留的参数知识。

任务说明、示范答案和查询对象承担不同职责。可以给模型指令“根据观察记录输出病害复查等级”，提供“叶片无斑点、边缘完整$\rightarrow$低”“叶片有扩散斑点$\rightarrow$高”两条教学示范，再输入一个没有答案的新记录。若把指令改为“输出物种名称”，原先示范就不再提供相同目标的证据。使用真实图片时还需要相应的图像输入接口；只有文本输入的语言模型不能直接检查未编码的图像。输入序列必须明确边界，查询答案只出现在评测端，不能附在其提示中。

任务状态可以通过几种不同计算形成。显式\term{任务后验推断}（task posterior inference）给潜在任务变量$u$规定先验和似然，计算
\begin{equation}
  p(y\mid\bm x,S_T,D_T)
  =\int p(y\mid\bm x,u)\,
          p(u\mid S_T,D_T)\,\mathrm du.
  \label{eq:ket-transfer-task-posterior}
\end{equation}
\begin{samepage}
这里的$u$可以表示规则、任务身份或参数，取离散值时积分改为求和。例如，病害案例中的一个候选规则可以是“出现扩散斑点就判为高复查等级”。式~\eqref{eq:ket-transfer-task-posterior}假定$u$足以确定标签机制，且查询输入不提供超出支持证据的任务身份信息；若任务还决定输入分布，就应把后验写成$p(u\mid\bm x,S_T,D_T)$，让查询输入也参与推断。

\end{samepage}

记忆读取从外部支持记录中找出相关条目，生成读取向量或选中的示范；这些条目不必被压成一个概率分布。上下文条件化直接以任务描述和示例作为网络输入，得到预测或像式~\eqref{eq:ket-transfer-cnp}那样的摘要。任务后验、记忆读取与上下文条件化分别强调状态的概率含义、证据的访问方式和预测对输入证据的依赖，可以同时成立。例如，网络激活可以编码后验参数，注意力也可以读取外部样本记录。讨论更新、撤销与使用成本时，还须说明实际保存的是后验参数、样本记录还是中间激活，以及它们保留多久。

Xie等在混合隐马尔可夫模型生成文档的设定中，分析语言建模怎样通过推断文档级潜在概念形成上下文能力。其论证假定模型能拟合预训练分布，并对提示与预训练的分布差异、潜在概念的可区分性作出限制；在这些条件下，式~\eqref{eq:ket-transfer-task-posterior}提供了可分析的解释。它没有证明任意实际语言模型都保存了正确任务先验或执行精确后验积分\scite{xie2022bayesianicl}

上下文计算也可能实现与参数更新等效的变换。对线性回归$\hat y=\bm w^\top\bm z$和平方损失$L(\bm w)=\frac1{2n}\sum_i(\bm w^\top\bm z_i-y_i)^2$，一次梯度更新后的查询预测为
\begin{equation}
  \hat y(\bm z)
  =\bm w^{(0)\top}\bm z
   -\frac{\alpha}{n}\sum_{i=1}^{n}
     (\bm w^{(0)\top}\bm z_i-y_i)\,\bm z_i^\top\bm z.
  \label{eq:ket-transfer-context-gd}
\end{equation}
令支持残差$r_i=\bm w^{(0)\top}\bm z_i-y_i$，取查询$\bm q=\bm z$、键$\bm k_i=\bm z_i$、标量值$v_i=-\alpha r_i/n$。不做softmax，也不除以权重和，直接读取$\sum_i(\bm q^\top\bm k_i)v_i$，所得便是式~\eqref{eq:ket-transfer-context-gd}中加在初始预测上的修正项。因此，这种未归一化点积读取可以在前向计算中得到同一预测，无需把$\bm w^{(1)}$写进持久模型权重。von Oswald等给出了实现这种线性注意力的权重构造，并研究了简单回归任务上训练得到的注意力模型；其构造依赖特定数据编码。这说明这里的代数对应，把它推广到任意语言任务或任意Transformer则没有相应依据\scite{ket-vonoswald2023gradient}

\clearpage
是否真正利用了支持证据，需要控制实验而不能只看输出是否变化。Min等在一组分类与多项选择任务上发现，示范的标签空间、输入分布与格式可能比每个示范标签的正确性更能解释收益；这个观察不表示任意新映射都能忽略正确标签\scite{ket-min2022demonstrations}表~\ref{tab:ket-transfer-context-controls}给出花卉任务可采用的对照，所有条件应保持查询集和基本证据预算一致。

\begin{table*}[tbp]
  \centering
  \small
  \begin{tabularx}{\linewidth}{p{30mm}XX}
    \toprule
    受控变化 & 保持不变的对象 & 希望区分的解释 \\
    \midrule
    将示范标签随机替换 & 支持输入、标签词表、序列长度 & 是否依赖正确输入答案对应，还是主要利用格式与标签范围 \\
    将类别名改为新符号并给出一致映射 & 输入和正确对应关系 & 是否学会当前映射，还是沿用预训练中的类别名称习惯 \\
    交换示范顺序 & 全部输入答案对与查询 & 顺序或邻近位置是否主导结果 \\
    仅改变格式和分隔符 & 任务语义与所有示范内容 & 是否对表示接口敏感，而不是知识不足 \\
    更换任务指令或移除示范 & 查询输入与允许信息来源 & 任务说明和示范各自贡献什么，二者冲突时依赖谁 \\
    \bottomrule
  \end{tabularx}
  \caption{上下文证据使用的控制实验。标签替换、改名和指令变化应在单独条件下进行；合成映射为教学设计，不预设模型能够完成。提示选择使用验证任务，最终查询答案不参与选择。}
  \label{tab:ket-transfer-context-controls}
\end{table*}

示范检索可以按查询相似性选择支持记录，但检索库必须只含允许信息，不能把查询本身或同源答案取回来。序列长度增加也增加编码费用；因果语言模型可以缓存固定前缀的键值状态，但缓存仍占存储，示范改变时还需要重算。允许跨请求保留记忆时，应注明任务切换和状态清空规则，否则一次测试任务的证据可能无意泄漏到另一任务。强化学习中的任务推断还涉及主动选择动作获得反馈，其探索和回报问题属于第~\ref{part:reinforcement-learning}部分的内容范围，不能由被动示范读取直接替代。

\subsection{模型参数的调整}
\label{sec:ket-transfer-parameter-adaptation}

目标证据可以通过优化形成持久任务参数。描述参数操作至少需要两个维度：来源已有参数更新多少，以及是否创建任务专有参数。前者可取“不更新、更新一部分、全部更新”，后者可取“不新增、新增”；两轴可以任意组合。例如，只训练新头是“已有参数不更新+新增”，全量微调新任务头是“已有参数全量更新+新增”，LoRA配合解冻归一化量则是“已有参数部分更新+新增”。因此，下面三个小节是讲解重点，不是互斥分类。

\begin{table}[htbp]
  \centering
  \small
  \begin{tabularx}{\linewidth}{p{28mm}XX}
    \toprule
    已有参数更新范围 & 不新增任务参数 & 新增任务参数 \\
    \midrule
    不更新 & 直接使用或上下文条件化已有模型 & 新建头、适配器、LoRA、软提示 \\
    部分更新 & 调整已有头、偏置、归一化量或部分层 & 新建头并解冻高层；LoRA并调整已有归一化量 \\
    全部更新 & 输出语义兼容时全量微调原模型 & 新任务头与全部来源参数共同微调 \\
    \bottomrule
  \end{tabularx}
  \caption{目标端参数操作的两个正交维度。单个方案由一格或跨格组合描述；“新增参数少”不推出已有参数冻结，也不推出反向传播路径短。}
  \label{tab:ket-transfer-parameter-axes}
\end{table}

更新范围与准备方式又是两个维度；普通预训练、MAML起点、层次先验和学习更新器都可以与表~\ref{tab:ket-transfer-parameter-axes}中的不同操作组合。使用学习更新器时，按式~\eqref{eq:ket-transfer-learned-update}初始化当前任务状态，再在选定参数集合上执行支持反馈更新。所有方案都应从同一目标支持信息出发，按独立验证确定步长与停止条件，最终测试标签不得参与选择。

\subsubsection{全量参数的调整}
\label{sec:ket-transfer-full-finetuning}

\term{全量微调}（full fine-tuning）允许来源表示与预测部分共同改变。若输出语义兼容，可直接以来源参数$\bm{\theta}_S$为$\bm{\theta}_T^{(0)}$；若目标需要新预测头，则先按新语义建立头，再与来源骨干共同训练。后者包含“全量调整来源参数”和“学习新建参数”两项操作，不能把新头说成从来源继承。

以新物种分类为例，设目标概率为$p_{\bm{\theta}_T}(y\mid\bm x)$，可优化
\begin{equation}
  L_T(\bm{\theta}_T)
   =-\frac1{|S_T|}\sum_{(\bm x_i,y_i)\in S_T}
       \log p_{\bm{\theta}_T}(y_i\mid\bm x_i)
       +\lambda\Omega(\bm{\theta}_T),
  \label{eq:ket-transfer-finetuning-loss}
\end{equation}
其中$\Omega$可以是预先指定的正则项或来自先验的约束。以允许的训练小批次执行$T$步更新，保存验证表现所选的检查点；若没有目标验证标签，应使用元验证任务上预先确定的固定配置，不能拿最终查询集充当验证集。分类改为检测时须相应改变损失，式~\eqref{eq:ket-transfer-finetuning-loss}的单标签交叉熵不能直接评价坐标或对象匹配。

与从头训练比较时，目标结构、支持数据、增强与测试条件应保持一致，分别报告来源初始化和随机初始化在相同预算下的适应曲线。还可以在充分训练条件下比较最终效果，但这回答的已不是固定少量更新预算的问题。全量微调自由度大，能改变特征提取，却也可能在极少样本下拟合偶然背景；是否过拟合由独立样本表现判断，不能仅由参数多作结论。

若来源参数数目为$P$，全量训练需要为这些参数计算和保存梯度，Adam等优化器还保存相应状态，并需处理中间激活。每个任务单独保存完整副本时，任务存储也约为$P$个参数。目标端更新步数不多，并不保证这一开销比构建参照更小。来源能力是否改变则是另一个评价维度：目标风险下降不能证明旧任务仍然可用。

同时发生拍摄条件或观测空间变化时，可以组合第~\ref{sec:ket-update-weighting}节的分布校正与式~\eqref{eq:ket-update-crossmodal}的配对约束，但应核对它们需要的无标签目标批次、配对观测或旧数据是否被允许。若还要求保留旧能力，就把第~\ref{sec:ket-update-retention-validation}节的旧任务验证与相应约束一并加入维护流程；本节的目标是形成当前任务能力，而不是用一次目标测试替代长期检查。

\subsubsection{已有部分参数的调整}
\label{sec:ket-transfer-existing-parameters}

来源模型某些部分已经适用时，可以只允许部分已有参数改变。将$\bm{\theta}$分为冻结部分$\bm{\theta}_F$与可训练部分$\bm{\theta}_A$，目标训练为
\[
  \min_{\bm{\theta}_A}
  L_T(\bm{\theta}_F=\bm{\theta}_{S,F},\bm{\theta}_A).
\]
只调已有预测头、只调骨干最高几层，或只调既有偏置和归一化参数，都属于这一范围。这里的“已有”指参数在来源模型中已经存在，不取决于可训练参数占比，也不取决于部署时是否为它另存一个副本。

比较时可固定物种集合和输出语义，例如目标仍区分来源已有的月季与鸢尾，仅拍摄背景改变。只调原有分类头能够改变固定特征空间中的决策边界；开放最高层还可以重新组合高层特征；全量调整允许更早的视觉特征变化。若来源编码器把叶片细斑压缩掉，任何只读其最终表示的头都不能补回细斑。选择更新层不是无条件的“越深越好”，而要看新目标差异出现在哪一层，以及现有目标证据能约束多少自由度。

冻结不等于不计算。只有顶部头可训练、骨干与预处理确定时，可以预先缓存支持特征，后续头部更新不必重复运行骨干；在线查询仍需骨干前向。若可训练模块在骨干中间，损失对该模块的梯度还要经过后续冻结层传播。冻结层通常可以不计算权重梯度，却仍需要计算输入梯度，并可能保留相关激活。训练显存与推理延迟因此不会按可训练参数比例自动缩小。

ANIL提供了采用元学习起点的部分更新例子：其目标端只适应头，准备阶段的骨干却经过跨任务训练。若目标类别数改变，或新的病害任务使原头的输出含义不再适用，而重新创建$\bm W_T,\bm b_T$，这些参数属于附加任务参数。复制、初始化、替换和冻结的具体操作比“只训练头”这个简称更能决定分类。新建头同时解冻骨干末层，则是附加参数学习与已有部分更新的组合。

\subsubsection{附加任务参数的学习}
\label{sec:ket-transfer-added-parameters}

最直接的附加参数是新建预测头。冻结$h_{\bm{\theta}_S}$，在$d$维表示上创建$N$类线性头，其权重$\bm W_T\in\mathbb R^{N\times d}$、偏置$\bm b_T\in\mathbb R^N$：
\begin{equation}
  \bm p(\bm x)
   =\operatorname{softmax}(\bm W_T h_{\bm{\theta}_S}(\bm x)+\bm b_T).
  \label{eq:ket-transfer-new-head}
\end{equation}
向量$\bm p(\bm x)$的第$c$项是类别$c$的预测概率，目标端训练$N(d+1)$个新参数。以交叉熵拟合时，一个样本对权重和偏置的梯度分别为$(\bm p-\bm e_y)\bm z^\top$与$\bm p-\bm e_y$，其中$\bm e_y$为标签独热向量。原型分类虽然可在平方距离下改写成受约束的线性打分，却没有独立优化这些权重，二者的参数形成过程不同。

只加输出头仍受固定表示限制。若希望改变网络中间计算而不修改底座权重，可以加入小型任务模块。\term{适配器}（adapter）通常在某个隐藏表示$\bm h\in\mathbb R^d$上加入瓶颈残差：
\begin{equation}
  \bm h'=\bm h+
     \bm W_{\mathrm{up}}\,
       \sigma(\bm W_{\mathrm{down}}\bm h+\bm b_{\mathrm{down}})
       +\bm b_{\mathrm{up}},
  \label{eq:ket-transfer-adapter}
\end{equation}
其中$\bm W_{\mathrm{down}}\in\mathbb R^{r\times d}$、$\bm W_{\mathrm{up}}\in\mathbb R^{d\times r}$，瓶颈宽度$r\ll d$。每个模块有$2dr+r+d$个参数；初始化使增量接近零，可以让训练开始时的网络接近原底座。新增参数通过目标损失训练，部署时保留底座与该任务模块，前向计算增加降维、非线性和升维。每个输入位置的两次矩阵向量乘法约需$2dr$次乘加，另有激活、偏置和残差相加；插入位置数与序列长度都会放大这项费用。含非线性的适配器通常不能像线性权重增量那样直接合并成原来的一张矩阵。

Houlsby等在Transformer每层的注意力输出投影之后和前馈子层之后各插入一个串行适配器，位置在对应的原残差相加与归一化之前；原残差连接和适配器内部残差不是同一条连接。其配置还调整归一化与任务输出部分，因此核算原方法时不能只数瓶颈矩阵。若工程实现复制了任务专有归一化量，应列出复制的状态；按本章的参数来源分类，直接调整底座已有归一化量仍属于已有部分参数更新\scite{ket-houlsby2019-adapters}

\term{低秩适应}（Low-Rank Adaptation, LoRA）不在一条串行路径上加入非线性，而是为选定线性权重添加低秩增量。设$\bm W_0\in\mathbb R^{d_{\mathrm{out}}\times d_{\mathrm{in}}}$冻结，新增$\bm A\in\mathbb R^{r\times d_{\mathrm{in}}}$和$\bm B\in\mathbb R^{d_{\mathrm{out}}\times r}$，前向计算变为
\begin{equation}
  \bm h'=\bm W_0\bm h+s\bm B\bm A\bm h,
  \label{eq:ket-transfer-lora}
\end{equation}
其中$s$为规定的缩放系数，原形式常取$s=\alpha_{\mathrm L}/r$。新增参数为$r(d_{\mathrm{in}}+d_{\mathrm{out}})$，且$\operatorname{rank}(\bm B\bm A)\leq r$。低秩是对这次权重增量施加的结构约束，不是对目标知识的语义结构作出发现，也不意味着底座权重本身变成低秩\scite{ket-hu2021lora}

LoRA常将$\bm A$随机初始化、$\bm B$置零，使初始增量为零而仍能开始学习。若令$\bm\delta=\partial L/\partial\bm h'$，单个输入的梯度为
\[
  \frac{\partial L}{\partial\bm B}
       =s\bm\delta(\bm A\bm h)^\top,\qquad
  \frac{\partial L}{\partial\bm A}
       =s\bm B^\top\bm\delta\,\bm h^\top.
\]
在$\bm B=0$的第一步，$\bm A$梯度为零，但$\bm B$通常有非零梯度；若两者都置零，则两个梯度都为零，增量无法从这个状态启动。这说明参数少并没有取消初始化设计。

当任务固定且数值格式兼容时，可以预先计算$\bm W=\bm W_0+s\bm B\bm A$，把增量合并进同形状权重，预测时不再执行额外分支；合并本身约需$O(d_{\mathrm{out}}d_{\mathrm{in}}r)$计算。多任务动态切换时，可以保留共享$\bm W_0$及各任务的$\bm A,\bm B$，但未合并分支每次输入增加$O(r(d_{\mathrm{in}}+d_{\mathrm{out}}))$计算。量化、并发多任务和切换频率会影响哪种部署方式合适，不能把“可合并”直接当作所有服务形态均无额外费用。

\term{软提示}（soft prompt）把可学习的连续向量放到输入端。若原输入嵌入矩阵为$\bm E\in\mathbb R^{\ell\times d}$，新增$\bm P_T\in\mathbb R^{m\times d}$，将按行拼接的$[\bm P_T;\bm E]$交给冻结模型，优化目标答案的负对数似然。训练对象是$md$个提示参数；提示通常不对应可直接阅读的词语。Lester等的prompt tuning主要在输入嵌入处加入这组向量，与在各层注入键值前缀的其他设计不能混为同一种参数计数\scite{ket-lester2021prompt}

软提示需要训练信号和梯度传播。虽然底座参数冻结，为计算输入提示的梯度，仍需通过模型各层反向传播。部署时保存$\bm P_T$和底座，输入长度由$\ell$变为$\ell+m$；在稠密自注意力中，仅注意力分数项就从$O(\ell^2d)$变为$O((\ell+m)^2d)$，增量为$O((2\ell m+m^2)d)$。它既占用上下文长度，也可能增加推理缓存，不能因为$md$很小就视为没有逐次费用。把文本示例直接放进上下文没有这一步参数训练，二者的证据使用和持久状态不同。

图~\ref{fig:ket-transfer-added-locations}比较新增参数进入前向计算的位置。输出头改变如何读取已有表示，适配器改变中间隐藏计算，LoRA改变某个线性映射的增量，软提示改变输入条件；这些位置约束了它们能表达的调整，但仅凭位置无法预言相同任务上的效果。

\begin{figure*}[htbp]
  \centering
  % Four forward-computation mechanisms; all dimensions are in millimetres.
% No external libraries beyond arrows.meta, positioning and calc are needed.
\begin{tikzpicture}[
  x=1mm,y=1mm,
  font=\figurefont\footnotesize,
  text=BookInk,
  >={Latex[length=2mm,width=1.4mm]},
  frozen/.style={draw=BookMuted,fill=BookPaper,line width=0.8pt,
    text width=30mm,minimum height=11mm,align=center,inner sep=2mm},
  train/.style={draw=BookTeal,fill=FigureModel!10,line width=0.8pt,
    text width=30mm,minimum height=11mm,align=center,inner sep=2mm},
  data/.style={draw=BookBlue,fill=FigureInput!10,line width=0.8pt,
    text width=21mm,minimum height=11mm,align=center,inner sep=2mm},
  add/.style={circle,draw=BookInk,fill=white,line width=0.8pt,
    minimum size=6mm,inner sep=0pt},
  flow/.style={draw=BookInk,line width=0.9pt,-{Latex},rounded corners=1mm},
  wire/.style={draw=BookInk,line width=0.9pt,rounded corners=1mm},
  heading/.style={anchor=west,font=\figurefont\small}
]
  \node[anchor=west,text=BookMuted] at (0,24)
    {共同条件：来源已有参数冻结；彩色模块：新增并训练的任务参数};
  \node[heading] at (0,13) {(a) 新建输出头：读取固定表示};
  \node[data] (head-input) at (13,0) {目标输入$\bm x$};
  \node[frozen] (backbone) at (61,0) {冻结底座\\$h_{\bm{\theta}_S}$};
  \node[train] (head) at (113,0) {训练新建头\\$\bm W_T,\bm b_T$};
  \node (head-output) at (157,0) {$\hat y$};
  \draw[flow] (head-input.east) -- (backbone.west);
  \draw[flow] (backbone.east) -- (head.west);
  \draw[flow] (head.east) -- (head-output.west);

  \node[heading] at (0,-22) {(b) 适配器：隐藏表示上的瓶颈残差};
  \node (adapter-input) at (10,-43) {$\bm h$};
  \coordinate (adapter-split) at (24,-43);
  \node[train] (down) at (52,-43) {训练降维\\再经非线性$\sigma$};
  \node[train] (up) at (103,-43) {训练升维\\$\bm W_{\mathrm{up}},\bm b_{\mathrm{up}}$};
  \node[add] (adapter-add) at (140,-43) {$+$};
  \node (adapter-output) at (158,-43) {$\bm h'$};
  \draw[wire] (adapter-input.east) -- (adapter-split);
  \draw[flow] (adapter-split) -- (down.west);
  \draw[flow] (down.east) -- (up.west);
  \draw[flow] (up.east) -- (adapter-add.west);
  \draw[flow] (adapter-split) -- (24,-31) -- (140,-31) -- (adapter-add.north);
  \draw[flow] (adapter-add.east) -- (adapter-output.west);

  \node[heading] at (0,-66) {(c) LoRA：线性映射上的并行低秩增量};
  \node (lora-input) at (10,-87) {$\bm h$};
  \coordinate (lora-split) at (24,-87);
  \node[frozen] (base-weight) at (78,-79) {冻结$\bm W_0$};
  \node[train] (low-rank) at (78,-99) {训练$s\bm B\bm A$};
  \node[add] (lora-add) at (140,-87) {$+$};
  \node (lora-output) at (158,-87) {$\bm h'$};
  \draw[wire] (lora-input.east) -- (lora-split);
  \draw[flow] (lora-split) |- (base-weight.west);
  \draw[flow] (lora-split) |- (low-rank.west);
  \draw[flow] (base-weight.east) -| (lora-add.north);
  \draw[flow] (low-rank.east) -| (lora-add.south);
  \draw[flow] (lora-add.east) -- (lora-output.west);

  \node[heading] at (0,-119) {(d) 软提示：输入序列中的连续任务条件};
  \node[train,text width=22mm] (prompt) at (15,-136) {训练提示\\$\bm P_T$};
  \node[data,text width=22mm] (tokens) at (15,-155) {输入嵌入\\$\bm E$};
  \node[data,text width=27mm] (concat) at (66,-145.5) {按序列位置拼接\\$[\bm P_T;\bm E]$};
  \node[frozen] (prompt-model) at (116,-145.5) {冻结模型\\输入长度$\ell+m$};
  \node (prompt-output) at (160,-145.5) {$\hat y$};
  \draw[flow] (prompt.east) -- (38,-136) |- (concat.west);
  \draw[flow] (tokens.east) -- (38,-155) |- (concat.west);
  \draw[flow] (concat.east) -- (prompt-model.west);
  \draw[flow] (prompt-model.east) -- (prompt-output.west);
\end{tikzpicture}

  \caption{附加任务参数的四种进入位置。实线为前向数据流，标有“训练”的模块包含新参数，标有“冻结”的模块保留来源权重；圆形加号表示同维向量相加。适配器面板只展开其内部瓶颈残差，实际插入位置见正文；LoRA面板的两条支路读取同一输入。软提示面板按序列位置拼接输入，不把连续提示当作已经给出的文本示范。}
  \label{fig:ket-transfer-added-locations}
\end{figure*}
\FloatBarrier

\begin{example}[固定底座下的参数与计算核算]
  \label{ex:ket-transfer-parameter-counts}
  设一个教学底座有12个Transformer层、隐藏宽度$d=768$，新任务有5类。这里只核算明确指定的附加部分，不宣称它们在这个任务上效果相同。新建线性头需要$5(768+1)=3845$个参数。

  若每层插入两个宽度$r=16$的适配器，每个为$2\cdot768\cdot16+16+768=25360$个参数，24个共608640个；再加新头为612485个。若另行训练每层两个已有LayerNorm的缩放与偏置，还需计入$12\cdot2\cdot2\cdot768=36864$个可训练量，总数变为649349。后者属于已有部分更新与附加参数的组合。

  若每层仅在查询和值两个$768\times768$矩阵上加秩$r=8$的LoRA，每矩阵增加$8(768+768)=12288$个参数，共294912个，加新头为298757个。每个未合并矩阵分支、每个输入位置需约12288次乘加；若合并为普通权重，则这部分逐位置分支计算消失，但合并与任务切换费用仍在。

  20个输入软提示需要$20\cdot768=15360$个参数。若它复用现成生成式输出接口，这就是提示参数数目；若同时新建相同的五分类头，则合计19205。这里仅核算单个注意力头的分数矩阵，不据此预测整个模型的推理时间：原输入为128个位置时，分数矩阵从$128^2=16384$个元素变为$148^2=21904$个，增加5520个，约33.69\%。
\end{example}

少量任务参数也可以在适应前经过跨任务准备，但要分清哪些模块在何时改变。Meta-Adapters包含冻结预训练底座$\bm{\theta}_S$、目标端可调的普通适配器$\bm\phi$，以及放在普通适配器前后的元适配器$\bm\omega$。历史任务的内层更新普通适配器，外层根据独立查询损失训练元适配器，同时学习普通适配器起点和逐层内层步长。一步形式可以写成
\begin{align}
  \bm\phi_k'
    &=\bm\phi_0-\bm\alpha\odot
       \nabla_{\bm\phi}L_{S_k}
         (\bm{\theta}_S,\bm\phi_0,\bm\omega), \nonumber\\
  \min_{\bm\omega,\bm\phi_0,\bm\alpha}\quad
  &\frac1M\sum_{k=1}^{M}
    L_{Q_k}(\bm{\theta}_S,\bm\phi_k',\bm\omega).
  \label{eq:ket-transfer-meta-adapters}
\end{align}
其中逐层步长在式中广播到对应参数坐标。到达新任务时，底座和元适配器都冻结，只从准备好的$\bm\phi_0$出发调整普通适配器；冻结的元适配器仍参与前向计算，也影响梯度经过它时的传播。因而该方法不仅学习一个小模块的起点，还学习“怎样调节底座以便小模块容易适应”。共享保存的$\bm\omega$和逐任务保存的适配器参数应分开计数，任务头也须纳入其实际可训练集合\scite{ket-bansal2022metaadapters}

对同一个新物种任务，表~\ref{tab:ket-transfer-operation-cost}归纳不同目标端操作。令$n$为支持样本数，$N$为类别数，$d$为表示维数，$T$为更新次数，$b$为训练小批次大小，$F$为单个输入的底座前向费用；$B_{\mathcal A}$表示对所选可训练集合$\mathcal A$完成一次样本反向与更新的费用。序列模型的前向费用另写作$F(\ell)$，因其随上下文长度变化。表中复杂度只描述指定实现的主要项，更新器、检索索引和实际硬件费用需单独实测。

\begin{table*}[htbp]
  \centering
  \small
  \setlength{\tabcolsep}{4pt}
  \renewcommand{\arraystretch}{1.2}
  \begin{tabularx}{\linewidth}{@{}>{\raggedright\arraybackslash}p{24mm}
      >{\raggedright\arraybackslash}p{36mm}
      >{\raggedright\arraybackslash}X
      >{\raggedright\arraybackslash}X@{}}
    \toprule
    操作 & 保存的任务状态 & 一次性适应费用 & 逐次预测费用 \\
    \midrule
    语义参照 & 属性表或约$Nd$个向量参照分量 & 名称、属性或描述的整理与编码；不必优化模型 & 图像编码与$N$类比较；向量形式约$F+O(Nd)$ \\
    \addlinespace
    支持样本或原型 & 逐例约$nd$个表示分量，原型约$Nd$个；另存标签 & 提取$n$个带标签样本的特征约$nF$，原型汇总另需$O(nd)$ & 逐例约$F+O(nd)$，原型约$F+O(Nd)$ \\
    \addlinespace
    上下文条件化 & 许可的说明、示范与历史，或其摘要、键值缓存 & 检索与组织证据，计算摘要或固定前缀 & 依接口为$F(\ell)$；长上下文与生成另计 \\
    \addlinespace
    预测头拟合 & 已有头或新建头，约$N(d+1)$个参数 & 可缓存固定特征时约$nF+O(TbNd)$ & 底座前向与头部计算，约$F+O(Nd)$ \\
    \addlinespace
    全量微调 & 完整任务权重或等价差量，含必要的新头 & 约$Tb(F+B_{\mathcal A})$；$\mathcal A$含全部来源参数及新头 & 适应后完整模型前向 \\
    \addlinespace
    附加任务参数 & 适配器、低秩增量或软提示参数，以及底座版本 & 前向、相关反向与新参数更新；LoRA合并和切换另计 & 底座前向加模块或输入位置开销；LoRA合并后无额外分支 \\
    \bottomrule
  \end{tabularx}
  \caption{相同允许信息下的适应操作与成本。方案不必用尽全部证据，也可占多行；新建头同时属于头部拟合与附加参数。参数优化使用任务监督，配置由独立验证确定。支持比较按固定独立表示计，完整上下文编码器还需额外计算；LoRA合并要求任务固定且数值格式兼容。基础编码器、语义模型和元训练的准备费用另按式~\eqref{eq:ket-reuse-amortization}核算，不混入表内两类目标端费用。}
  \label{tab:ket-transfer-operation-cost}
\end{table*}

表中的参数数量、标签需求、训练时间与预测代价是四个不同指标。软提示可以参数少而反向路径长；原型可以没有梯度更新而使用全部支持标签；全量微调可能目标端代价较高，却在大量重复预测时摊薄一次性适应费用。费用比较必须使用同一允许证据和目标标准，也应记录示例检索、提示组织、上下文长度、参数更新、推理延迟与任务状态存储。

若目标明显超出历史覆盖，例如来源只做整图物种识别，而目标病害识别依赖微小叶斑，则应检查当前表示和观测是否保留必要信息。适应曲线在增加样本和更新次数后仍无改善时，可以补充有判别力的观测与标注、扩大可训练范围，或回退到普通目标训练；不能仅因方案被称为元学习而坚持固定的一步适应。检验收益应画出按标注数、更新次数或真实计算预算变化的目标表现，而不是只挑一个有利预算点。目标能力形成后若继续承担历史与后续任务，其版本、记忆和长期成本须纳入第~\ref{sec:ket-update-maintenance}节的持续维护。

\section{本章小结}
\label{sec:ket-transfer-summary}

已有花卉识别知识能否帮助一个新植物园，取决于它与新目标之间存在什么可用联系。新物种要求补足类别语义和区分证据，病害判断要求表示保留病害相关信息，定位则要求位置与对象级预测结构。识别为未知只解决现有类别目录不足时的退出问题，既没有自动区分新类别，也没有建立可验证的新类别能力。

来源经验与目标证据可以通过不同方式相互补足：在表示空间建立参照，用语义描述连接未见类别，让合适起点或先验约束少样本估计，使用准备好的更新器，或把支持证据读成任务上下文。元学习的特殊之处是准备目标包含了适应后的表现，而不是它必然更新更多或更少参数。已有部分参数与新建参数的边界由参数来源决定；适配器、LoRA和软提示节省的资源也不相同，必须按前向位置、反向路径、状态存储和部署方式核算。

有说服力的迁移证据需要同时保留任务内与任务间的数据边界，并在固定信息和预算下观察独立目标的适应曲线。少样本、少参数或少量更新各自说明一个局部条件，不能单独证明整体学习更高效。准备费用应按实际使用任务摊销，并把失败尝试计入成本；只有完整账目与目标效果的比较才能说明是否值得准备。一次目标能力形成也不能替代长期维护已有能力的验证。
